% from main.tex
%%%%%%%%%%  scan idx 532  |  folio 519  %%%%%%%%%%
\section*{Exposé X. Formalisme des intersections sur les schémas algébriques propres}
\begin{center}
par O. Jussila\\
(avec un appendice par A. Grothendieck)
\end{center}

\section*{1. Compatibilité des filtrations avec la loi de composition}
\[
K^\circ(X)\times K_\bullet(X)\longrightarrow K_\bullet(X).
\]

\subsection*{1.1}
Considérons d'abord un schéma noethérien $X$ et le groupe $K_\bullet(X)$ des classes de $\mathcal O_X$-modules cohérents. Pour tout $\mathcal O_X$-module cohérent $F$ nous désignerons par $\cl_\bullet(F)$ sa classe dans $K_\bullet(X)$.

\textbf{Définition 1.1.1.} Pour tout entier $n\in\mathbb Z$ on désigne par $\Filt_n(X)$ le sous-groupe de $K_\bullet(X)$ engendré par les classes $\cl_\bullet(F)$ des $\mathcal O_X$-modules cohérents $F$ tels que
\[
\dim(\Supp(F))\leq n.
\]
On définit ainsi une filtration croissante $(\Filt_n(X))_{n\in\mathbb Z}$ de $K_\bullet(X)$ telle que $\Filt_n(X)=0$ pour $n<0$.

\textbf{Proposition 1.1.2.} Soient $X$ un schéma noethérien, $F$ un $\mathcal O_X$-module cohérent tel que $\dim(\Supp(F))\leq n$, et $X(n)$ ($n\in\mathbb Z$) le sous-ensemble $\{x\in X\,;\,\dim(\overline{\{x\}})=n\}$ de $X$. Pour tout $x\in X$ on désigne par $Z(x)$ le sous-schéma intègre fermé associé à $\overline{\{x\}}$. Alors la longueur du $\mathcal O_{X,x}$-module $F_x$ est finie pour tout $x\in X(n)$ et non-nulle pour un nombre fini des $x\in X(n)$ seulement, et on a
